\section{Related Tutorials and Where This One Differs}
\label{sec:related}

Three bodies of work meet in this tutorial, and it is useful to say up front what we
borrow from each and what we add.

\textit{Graph neural networks.} Surveys of GNNs~\cite{wu2021gnnsurvey} catalog
architectures (spectral~\cite{bruna2014spectral,defferrard2016chebnet},
convolutional~\cite{kipf2017gcn}, attentional~\cite{velickovic2018gat},
message-passing~\cite{gilmer2017mpnn}, expressive~\cite{xu2019gin}, and
inductive~\cite{hamilton2017graphsage}) and their applications. A parallel line studies
the depth and range limits of these models, including over-smoothing fixes such as
deep propagation~\cite{chen2020gcnii,gasteiger2019appnp} and the over-squashing and
bottleneck phenomena that bound how far a local operator can carry
information~\cite{alon2021bottleneck,topping2022oversquashing}. We do not add to that
catalog. Instead we adopt the spectral, operator-theoretic reading of
GNNs~\cite{shuman2013gsp} and push it to its conclusion: once a layer is written as a
function of the Laplacian, the classical and quantum cases sit on the same axis.
Diffusion-based learning~\cite{gasteiger2019diffusion} is the closest classical neighbor
of our view, and dynamic-graph learning~\cite{rossi2020tgn,kazemi2020dynamicsurvey}
supplies the time-varying setting.

\textit{Quantum machine learning.} Reviews of quantum machine
learning~\cite{biamonte2017qml,dunjko2018review} and variational quantum
algorithms~\cite{cerezo2021vqa} survey models~\cite{schuld2014quest}, feature
maps~\cite{schuld2019featurespace,havlicek2019supervised}, trainability and
power~\cite{abbas2021power}, and the obstacles of the NISQ
era~\cite{preskill2018nisq,mcclean2018barren}. Quantum graph models in particular,
including quantum graph neural networks~\cite{verdon2019qgnn}, equivariant quantum graph
circuits~\cite{mernyei2022equivariant}, quantum graph
kernels~\cite{henry2021quantumkernel}, and the broad quantum-walk
literature~\cite{farhi1998quantumwalk,kempe2003qwalk,childs2003exponential,childs2004quantumwalk,childs2009universal,venegas2012quantumwalks,portugal2013book},
supply the operators we use. These works establish what is possible; they rarely
isolate, on a controlled task, what the quantum operator changes relative to a
param-matched classical one. That isolation is our contribution.

\textit{Learning for non-terrestrial networks.} The networking literature surveys
satellite systems~\cite{kodheli2021ntn}, studies low-latency routing over LEO
constellations~\cite{handley2018delay}, and provides realistic
simulators~\cite{kassing2020hypatia}. Graph and learning-based methods are by now
established for network modeling and resource
management~\cite{rusek2020routenet,shen2021gnnradio,eisen2020optimal,mao2016resource}.
We borrow the dynamic-graph setting from this line of work but deliberately keep the
task generic and the simulator minimal, so that the tutorial's conclusions are about
operators rather than about one deployment.

\textit{What is new here.} The gap these literatures leave is a single, teachable
through-line from classical message passing to quantum walks, paired with an honest
methodology for deciding whether the quantum option helps. We fill it with the
propagation-operator view (Sections~\ref{sec:classical}--\ref{sec:quantum}), a built,
runnable model (Section~\ref{sec:buildit}), a reproducible and multi-seed case study
(Section~\ref{sec:casestudy}), and a pitfall catalogue with a checklist
(Section~\ref{sec:pitfalls}). The aim is not breadth of coverage but a reader who can
both build and judge.
